\SourceLeaf{249}%
\[
\begin{aligned}
4EEGGk&=E\mathord{.}\frac{\mathrm{d}E}{\mathrm{d}q}\mathord{.}\frac{\mathrm{d}G}{\mathrm{d}q}+E(\frac{\mathrm{d}G}{\mathrm{d}p})^2+G\mathord{.}\frac{\mathrm{d}E}{\mathrm{d}p}\mathord{.}\frac{\mathrm{d}G}{\mathrm{d}p}\\
&\quad+G(\frac{\mathrm{d}E}{\mathrm{d}q})^2-2EG(\frac{\mathrm{d}\mathrm{d}E}{\mathrm{d}q^2}+\frac{\mathrm{d}\mathrm{d}G}{\mathrm{d}p^2})
\end{aligned}
\]
et pro variatione anguli \(\theta\)
\[
\BareRoot{EG}\mathord{.}\mathrm{d}\theta=\tfrac{1}{2}\mathord{.}\frac{\mathrm{d}E}{\mathrm{d}q}\mathord{.}\mathrm{d}p-\tfrac{1}{2}\mathord{.}\frac{\mathrm{d}G}{\mathrm{d}p}\mathord{.}\mathrm{d}q
\]
\par\indent Inter varios casus, in quibus haec conditio orthogonalitatis valet, primarium locum tenet is, ubi lineae omnes alterutrius systematis, e. g. primi, sunt lineae brevissimae. Hic itaque pro valore constante ipsius \(q\), angulus \(\theta\) fit \(=0\), unde aequatio pro variatione anguli \(\theta\) modo tradita docet, fieri debere \(\frac{\mathrm{d}E}{\mathrm{d}q}=0\), sive coëfficientem \(E\) a \(q\) independentem, i. e. \(E\) esse debet vel constans vel functio solius \(p\). Simplicissimum erit, pro \(p\) adoptare longitudinem ipsam cuiusque lineae primi systematis, et quidem, quoties omnes lineae primi systematis in uno puncto concurrunt, ab hoc puncto numeratam, vel, si communis intersectio non adest, a qualibet linea secundi systematis. Quibus ita intellectis patet, \(p\) et \(q\) iam eadem denotare, quae in artt. 15, 16 per \(r\) et \(\varphi\) expresseramus, atque fieri \(E=1\). Ita duae formulae praecedentes iam transeunt in has:
\[
\begin{aligned}
4GGk&=(\frac{\mathrm{d}G}{\mathrm{d}p})^2-2G\frac{\mathrm{d}\mathrm{d}G}{\mathrm{d}p^2}\\
\BareRoot{G}\mathord{.}\mathrm{d}\theta&=-\tfrac{1}{2}\mathord{.}\frac{\mathrm{d}G}{\mathrm{d}p}\mathord{.}\mathrm{d}q
\end{aligned}
\]
vel statuendo \(\BareRoot{G}=m\),
\[
k=-\frac{1}{m}\mathord{.}\frac{\mathrm{d}\mathrm{d}m}{\mathrm{d}p^2},\qquad\mathrm{d}\theta=-\frac{\mathrm{d}m}{\mathrm{d}p}\mathord{.}\mathrm{d}q
\]
Generaliter loquendo \(m\) erit functio ipsarum \(p,q\) atque \(m\mathrm{d}q\) expressio elementi cuiusvis lineae secundi systematis. In casu speciali autem, ubi omnes lineae \(p\) ab eodem puncto proficiscuntur, manifesto pro \(p=0\) esse debet \(m=0\); porro si in hoc casu pro \(q\) adoptamus angulum ipsum, quem elementum primum cuiusvis lineae primi systematis facit cum elemento alicuius ex ipsis ad arbitrium electae, quum pro valore infinite parvo ipsius \(p\), elementum lineae secundi systematis (quae considerari potest tamquam circulus radio \(p\) descriptus), sit \(=p\mathrm{d}q\), erit pro valore infinite parvo ipsius \(p,m=p\), adeoque, pro \(p=0\) simul \(m=0\) et \(\frac{\mathrm{d}m}{\mathrm{d}p}=1\).
\Article{20}
Immoremur adhuc eidem suppositioni, puta \(p\) designare indefinite longitudinem lineae brevissimae a puncto determinato \(A\) ad punctum quodlibet super%
